%-------------------------------------------------------------------------------
%	SECTION TITLE
%-------------------------------------------------------------------------------
\cvsection{Projets}


%-------------------------------------------------------------------------------
%	CONTENT
%-------------------------------------------------------------------------------
\begin{cventries}

	%---------------------------------------------------------
	\cventry
	{Claude}
	{Jobs Assistant - Agent IA pour la veille de missions freelance}
	{HomeLab}
	{2026}
	{
		\begin{cvitems} % Description(s) of experience/contributions/knowledge
			\item {Critique de CV}
			\item {Recherche d'annonces basées sur un CV}
			\item {Rédaction de lettres de motivation en fonction de l'annonce}
			\item {Préparation aux entretiens}
		\end{cvitems}
	}

	%---------------------------------------------------------
	\cventry
	{Claude}
	{Shop Assistant - Agent IA pour l'interaction avec le site d'Intermarché}
	{HomeLab}
	{2026}
	{
		\begin{cvitems} % Description(s) of experience/contributions/knowledge
			\item {Création d'un outil CLI pour interagir avec le site d'Intermarché}
			\item {Utilisation de Chrome pour bypass les restrictions de Bots}
			\item {Consommation de l'outil CLI exposé en MCP pour faire des courses en se basant notamment sur les commandes précédentes}
		\end{cvitems}
	}

	%---------------------------------------------------------
	\cventry
	{Docker}
	{MCP Chrome (OpenSource)}
	{HomeLab}
	{2026}
	{
		\begin{cvitems} % Description(s) of experience/contributions/knowledge
			\item {Création d'un MCP Chrome}
			\item {Exposition de Chrome via le CDP}
			\item {Utilisation d'un GPU NVIDIA dans le conteneur pour avoir de l'\textbf{accélération matérielle} et éviter le bot flag.}
			\item {Utilisation de VNC pour un accès à distance et résoudre les Captcha manuellement}
		\end{cvitems}
	}

	%---------------------------------------------------------
	\cventry
	{Nix~~~·~~~NixOS~~~·~~~nix-darwin}
	{\href{https://gitlab.com/NicolasGuilloux/nixos-configuration}{NixOS Configuration (OpenSource)}}
	{HomeLab}
	{2020 - présent}
	{
		\begin{cvitems} % Description(s) of experience/contributions/knowledge
			\item {Configuration déclarative multi-machines : Linux (NixOS) et macOS (nix-darwin), gestion via \textbf{Nix Flakes} et Snowfall Lib}
			\item {Gestion des secrets chiffrés (\textbf{ragenix}), environments de développement (\textbf{devenv}), thématisation système (Stylix)}
			\item {Représentant de la communauté Linux auprès de Shadow, testeur Insider}
		\end{cvitems}
	}

	%---------------------------------------------------------
	\cventry
	{Nix~~~·~~~NixOS}
	{\href{https://github.com/NicolasGuilloux/shadow-nix}{Shadow Nix Package (OpenSource)}}
	{HomeLab}
	{2020-2021}
	{
		\begin{cvitems} % Description(s) of experience/contributions/knowledge
			\item {Packaging du client Shadow, un service de Cloud Gaming français}
			\item {Création d'un Live ISO et d'un environnement minimal}
		\end{cvitems}
	}

	%---------------------------------------------------------
	\cventry
	{Symfony}
	{\href{https://github.com/rich-id/test-suite}{Suite de test} (OpenSource)}
	{Rich-ID}
	{2020-2021}
	{
		\begin{cvitems} % Description(s) of experience/contributions/knowledge
			\item {Suite de tests pour améliorer et étendre les fonctionnalités de celle de Symfony}
			\item {Limite les effets de bords, ajoute des fonctions pratiques, limite le boilerplate}
		\end{cvitems}
	}

	%---------------------------------------------------------
	\cventry
	{Symfony}
	{\href{https://github.com/NicolasGuilloux/chaplean-dto-handler-bundle}{DTO Handler Bundle (OpenSource)}}
	{Chaplean}
	{2019}
	{
		\begin{cvitems} % Description(s) of experience/contributions/knowledge
			\item {\textbf{ParamConverter} pour automatiquement charger le contenu d'une requête dans un \textbf{DTO}, et le \textbf{valider}}
		\end{cvitems}
	}

	%---------------------------------------------------------
\end{cventries}
